\documentclass[11pt]{article}

\usepackage[margin=1in]{geometry}
\usepackage{amsmath, amssymb, amsthm}
\usepackage{hyperref}
\usepackage{enumitem}
\usepackage{microtype}

% Theorem environments
\newtheorem{theorem}{Theorem}[section]
\newtheorem{lemma}[theorem]{Lemma}
\newtheorem{corollary}[theorem]{Corollary}
\newtheorem{proposition}[theorem]{Proposition}
\theoremstyle{definition}
\newtheorem{definition}[theorem]{Definition}
\newtheorem{example}[theorem]{Example}
\theoremstyle{remark}
\newtheorem{remark}[theorem]{Remark}
\newtheorem{conjecture}[theorem]{Conjecture}

% Notation shorthands
\newcommand{\F}{\mathcal{F}}
\newcommand{\G}{\mathcal{G}}
\newcommand{\U}{\mathcal{U}}
\newcommand{\Pow}{\mathcal{P}}
\newcommand{\card}[1]{\#\{#1\}}

\title{Frankl's Union-Closed Sets Conjecture\\[0.4em]
\large Math 475 Individual Report}
\author{Vishesh Narayan}
\date{November 12, 2026}

\begin{document}
\maketitle

\tableofcontents
\newpage

% ---------------------------------------------------------------
\section{The Problem}
% ---------------------------------------------------------------

\subsection{Setup and Notation}

For a positive integer $n$, write $[n] := \{1, 2, \ldots, n\}$ and let
$\Pow(n) := \Pow([n])$ denote the power set of $[n]$.

\begin{definition}
A family $\F \subseteq \Pow(n)$ is \emph{union-closed} if for every
$A, B \in \F$ we also have $A \cup B \in \F$.
\end{definition}

\begin{definition}
A family $\F \subseteq \Pow(n)$ is \emph{nontrivial} if
$\bigcup_{F \in \F} F = [n]$, i.e., every element of the ground set
appears in at least one member of $\F$.
Note that any nontrivial union-closed family must contain $[n]$ itself.
\end{definition}

\begin{example}
Let $n = 3$ and consider
\[
  \F = \bigl\{\{1\},\, \{2\},\, \{1,2\},\, \{1,3\},\, \{2,3\},\,
              \{1,2,3\}\bigr\}.
\]
One checks that $\F$ is union-closed and nontrivial.
The element frequencies are: element $1$ appears in $4$ of $6$ sets,
element $2$ appears in $4$ of $6$, element $3$ appears in $3$ of $6$.
In particular, both $1$ and $2$ appear in at least half the sets.
\end{example}

\subsection{The Conjecture}

\begin{conjecture}[Frankl's Union-Closed Sets Conjecture, $\sim$1979]
\label{conj:func}
For every nontrivial, union-closed family $\F \subseteq \Pow(n)$
there exists an element $x \in [n]$ such that
\[
  \#\{F \in \F : x \in F\} \;\ge\; \frac{1}{2} \cdot \#\F.
\]
\end{conjecture}

Despite its elementary statement, Conjecture~\ref{conj:func} has been
open since the 1980s.  No proof or counterexample is known in general.

\subsection{Equivalent Formulations}

The conjecture admits several equivalent restatements that illuminate
different aspects of the problem.

\textbf{Intersection-closed dual.}
For a family $\F \subseteq \Pow(n)$, define the \emph{complement family}
$\G := \{[n] \setminus F : F \in \F\}$.
Then $\F$ is union-closed if and only if $\G$ is
\emph{intersection-closed} (i.e., $A, B \in \G \Rightarrow A \cap B \in \G$).
Under complementation, an element $x$ is in $F \in \F$ if and only if
$x \notin [n]\setminus F \in \G$.
So the conjecture is equivalent to:

\begin{conjecture}[Intersection-Closed Dual]
For every intersection-closed family $\G \subseteq \Pow(n)$ with
$\bigcap_{G\in\G} G = \emptyset$, there exists $x \in [n]$ with
\[
  \#\{G \in \G : x \in G\} \;\le\; \frac{1}{2}\cdot\#\G.
\]
\end{conjecture}

Neither form is obviously easier to work with, but the intersection-closed
version is natural in lattice theory.

\textbf{Averaging reformulation.}
An immediate thought is whether a \emph{random} element of $[n]$
(chosen uniformly) belongs to at least half the sets in expectation.
This is false in general: for $\F = \{\{1\},\{1,2\},\{1,3\},\{1,2,3\}\}$
the average membership is $10/4 \cdot \frac{1}{3} < 2 = \frac{1}{2}\#\F$.
The conjecture is weaker --- it only asserts the \emph{existence} of one
good element, not that the average element is good \cite{gowers2016}.

\textbf{Interior-operator formulation.}
For a union-closed family $\F$ containing $\emptyset$, define
\[
  \tau : \Pow(n) \to \F, \qquad \tau(X) := \bigcup\{F \in \F : F \subseteq X\}.
\]
This map satisfies $\tau(X) \subseteq X$ (exclusivity),
$X \subseteq Y \Rightarrow \tau(X)\subseteq\tau(Y)$ (monotonicity),
and $\tau(\tau(X)) = \tau(X)$ (idempotence), making it an
\emph{interior operator}.
The correspondence $\F \leftrightarrow \tau$ is a bijection between
union-closed families containing $\emptyset$ and interior operators on
$\Pow(n)$ \cite{nagel2023}.
This viewpoint organises $\Pow(n)$ into \emph{congruence classes}
$\mathcal{T}(F) = \tau^{-1}(F)$ for each $F \in \F$, each of which is
closed under intersection and ``interval-complete'' (if $A \subseteq X \subseteq F$
and $A \in \mathcal{T}(F)$, then $X \in \mathcal{T}(F)$).

% ---------------------------------------------------------------
\section{What Is Known}
% ---------------------------------------------------------------

\subsection{Easy cases}
\label{sec:easy}

\textbf{Singleton case (folklore).}
If $\{x\} \in \F$ for some $x$, then for every $F \in \F$ we have
$F \cup \{x\} \in \F$, so the map $F \mapsto F \cup \{x\}$ is an
injection from $\{F \in \F : x \notin F\}$ into $\{F \in \F : x \in F\}$.
Hence element $x$ appears in at least half the sets.

\textbf{Two-element ground set.}
If $n = 2$ the conjecture is trivial by inspection of all union-closed
families on $\{1,2\}$.

\textbf{Small families.}
The conjecture is verified for $\#\F \le 50$ by exhaustive methods
\cite{bruhn2015}.

\subsection{Density bounds}

Let $\F \subseteq \Pow(n)$ be a nontrivial union-closed family of size
$m = \#\F$.
The best known lower bounds on the frequency of the most common element are
\[
  \Omega\!\left(\max\!\left\{\sqrt{\frac{\log_2 n}{n}},\;
    \frac{\log_2 m}{n},\; \frac{1}{\log_2 m}\right\}\right) \cdot m
\]
due to Balla~\cite{balla2011}, Reimer~\cite{reimer2003}, and
W\'{o}jcik~\cite{wojcik1999} respectively (the third improving a
constant from Knill).
These bounds confirm the conjecture when $m$ is large relative to $n$,
but the general case remains open.

\subsection{A weakening via up-sets}

Nagel~\cite{nagel2023} proves the following weakening of the conjecture.
An \emph{up-set} is a family $\U \subseteq \Pow(n)$ such that
$A \in \U$ and $A \subseteq B$ imply $B \in \U$.

\begin{theorem}[Nagel 2023]
\label{thm:upset}
Let $\F \subseteq \Pow(n)$ be a nontrivial union-closed family.
There exists an up-set $\U \subseteq \Pow(n)$ with
$\#\U \le 2^{n-1}$ such that
\[
  \#(\F \cap \U) \;\ge\; \tfrac{1}{2} \cdot \#\U.
\]
\end{theorem}

This is weaker than the conjecture because for the conjecture we would
need $\U = [x,[n]] := \{X \subseteq [n] : x \in X\}$ for some specific
$x$; Theorem~\ref{thm:upset} allows any up-set of bounded size.
Nevertheless the result is nontrivial: it shows that $\F$ always has
a large ``dense'' region in the Boolean lattice.

% ---------------------------------------------------------------
\section{Approaches and Ideas Explored}
% ---------------------------------------------------------------

\subsection{The injection paradigm}

The most natural proof strategy generalises the singleton argument of
Section~\ref{sec:easy}.  Fix any element $x \in [n]$ and partition
\[
  \F^- := \{A \in \F : x \notin A\}, \qquad
  \F^+ := \{A \in \F : x \in A\}.
\]
The conjecture asks for $|\F^+| \ge |\F^-|$.
If we can construct an \emph{injective} map
$\varphi : \F^- \to \F^+$, then immediately $|\F^+| \ge |\F^-|$.

The natural template is to pick, for each $A \in \F^-$, some set
$B(A) \in \F^+$ and set $\varphi(A) := A \cup B(A)$.
Since $\F$ is union-closed, $A \cup B(A) \in \F$; since $x \in B(A)$,
we have $x \in A \cup B(A)$, so $\varphi(A) \in \F^+$.
Injectivity of $\varphi$ depends on the choice of $B(A)$:
different sets $A, A'$ must produce distinct images $A \cup B(A) \neq A' \cup B(A')$.

The \emph{singleton injection} (K1) takes $B(A) = \{x\}$ for every $A$,
valid when $\{x\} \in \F$.
Injectivity follows because $A \cup \{x\} = A' \cup \{x\}$ implies
$A = A'$ (both sets agree on $[n] \setminus \{x\}$, and neither contains $x$).

\textbf{Nagel's fiber lemma.}
Nagel \cite{nagel2023} observes that the singleton argument is a special case of
a broader principle.

\begin{lemma}[Nagel 2023, Lemma 2.4]
\label{lem:fiber}
Let $\F \subseteq \Pow(n)$ be nontrivial and union-closed, $x \in [n]$, and
let $A \in \F$ with $x \in A$.  Then
\[
  \#\{F \in \F : x \in F\} \;\ge\; \frac{1}{2^{\#A - 1} + 1} \cdot \#\F.
\]
\end{lemma}

\begin{proof}[Proof sketch]
Define $\varphi : \F^- \to \F^+$ by $\varphi(G) = G \cup A$.
For each image $F = G \cup A$, the preimage satisfies
$\varphi^{-1}(F) = \{G' \subseteq [n]\setminus\{x\} : G' \cap ([n]\setminus A) = F \cap ([n]\setminus A)\}$,
since $G \cup A = F \Leftrightarrow G \setminus A = F \setminus A$.
The set $G$ can only vary on $A \setminus \{x\}$, giving $|\varphi^{-1}(F)| \le 2^{\#A-1}$.
Hence $|\F^-| \le 2^{\#A-1} \cdot |\F^+|$, and rearranging
$|\F| = |\F^-| + |\F^+| \le (2^{\#A-1}+1)|\F^+|$ gives the lemma.
\end{proof}

For $A = \{x\}$ (a singleton) the bound gives $|\F^+| \ge \frac{1}{2}|\F|$,
recovering the folklore result.
For larger $A$ the bound degrades, but crucially \emph{any} set $A \in \F^+$
yields a nontrivial lower bound: even when no singleton is present we can
apply the lemma with the smallest available $A \in \F^+$.

\subsection{Nagel's frequency question}

\label{sec:freq}

Label elements so that
$f_1 := \#\{F \in \F : 1 \in F\} \ge f_2 \ge \cdots \ge f_n$.
Conjecture~\ref{conj:func} says $f_1 \ge \tfrac{1}{2}|\F|$.
What can be said about $f_k$ for $k > 1$?

An iterative argument shows the following.
If the conjecture holds, the family $\F' := \{F \setminus \{1\} : 1 \in F \in \F\}$
is again nontrivial and union-closed of size $|\F'| \ge \tfrac{1}{2}|\F|$,
so by the conjecture applied to $\F'$ there exists $x \in [n] \setminus \{1\}$ with
$\#\{F \in \F : x \in F\} \ge \tfrac{1}{4}|\F|$,
and since element $2$ is at least as frequent as $x$, we get $f_2 \ge \tfrac{1}{4}|\F|$.
Iterating: $f_k \ge 2^{-k}|\F|$.
This bound is \emph{wasteful}, since at each step we discard all sets not
containing the previously chosen element.

Nagel \cite{nagel2023} conjectures the tight bound:
\begin{equation}
  \label{eq:freqq}
  f_k \;\ge\; \frac{1}{2^{k-1}+1}\cdot|\F| \qquad \text{for all } k \in [n].
\end{equation}
For $k=1$ this is exactly Conjecture~\ref{conj:func}.
The bound \eqref{eq:freqq} is sharp: take
$\F = \Pow(k-1) \cup \{[n]\}$; then $|\F| = 2^{k-1}+1$, element $k$ appears
only in $\{[n]\}$, so $f_k/|\F| = 1/(2^{k-1}+1)$.

Nagel proves \eqref{eq:freqq} for $k = n$ and $k = n-1$ using
Lemma~\ref{lem:fiber}: for $k = n$, apply the lemma with $x = n$ and $A = [n]$;
for $k = n-1$, exploit separability to find a set $A \in \F^+_{n-1}$ with
$|A| \le n-1$ and apply the lemma.
The cases $k < n-1$ remain open and are equivalent in difficulty to
the original conjecture.

\subsection{The up-set weakening}

We now outline the proof of Theorem~\ref{thm:upset} following Nagel \cite{nagel2023}.
Recall that the conjecture is equivalent to: $|\F \cap [x,[n]]| \ge \frac{1}{2}|\F|$
for some $x$, where $[x,[n]] := \{X \subseteq [n] : x \in X\}$ is an up-set of size $2^{n-1}$.
Theorem~\ref{thm:upset} relaxes this by allowing \emph{any} up-set.

\textbf{Setup.}
Assume without loss of generality that $\emptyset \in \F$.
Define the \emph{interior operator} $\tau : \Pow(n) \to \F$ by
$\tau(X) := \bigcup\{F \in \F : F \subseteq X\}$
(the largest element of $\F$ contained in $X$).
This partitions $\Pow(n)$ into fibres $T(F) := \tau^{-1}(F)$ for each $F \in \F$,
with $\bigsqcup_{F \in \F} T(F) = \Pow(n)$.

\textbf{Key properties.}
Each fibre $T(F)$ is intersection-closed and interval-complete
(if $A \subseteq X \subseteq F$ and $A \in T(F)$, then $X \in T(F)$),
so $T(F)$ is ``topologically nice''.
Crucially: if $E \subsetneq F$ in $\F$, then $|T(F)| \le |T(E)|$ (Lemma 3.5 of \cite{nagel2023}).
Larger sets in $\F$ have smaller fibres.

\textbf{Proof of Theorem~\ref{thm:upset}.}
Order $\F = \{F_1, \ldots, F_m\}$ with $F_1 = [n]$ and $|T(F_1)| \le \cdots \le |T(F_m)|$.
(Such an ordering exists by the lemma above and union-closedness.)
Set $U := \bigcup_{i=1}^{\lceil m/2 \rceil} T(F_i)$.

\emph{$U$ is an up-set.}
Take $X \in T(F_i)$ with $i \le \lceil m/2 \rceil$ and $Y \supseteq X$.
By monotonicity of $\tau$, $\tau(Y) \supseteq \tau(X) = F_i$,
so $\tau(Y) = F_j$ for some $F_j \supseteq F_i$ in $\F$.
By the ordering, $j \le i \le \lceil m/2 \rceil$, hence $Y \in T(F_j) \subseteq U$.

\emph{$|\F \cap U| \ge \frac{1}{2}|\F|$.}
By construction, $\F \cap U = \{F_1, \ldots, F_{\lceil m/2 \rceil}\}$.

\emph{$|U| \le 2^{n-1}$.}
The fibre sizes $f_i := |T(F_i)|$ satisfy $f_1 = 1$ and $\sum_{i=1}^m f_i = 2^n$.
Since $f_1 \le \cdots \le f_m$, Lemma 3.9 of \cite{nagel2023}
(a monotone summation estimate) gives
$|U| = \sum_{i=1}^{\lceil m/2 \rceil} f_i \le \frac{1}{2}\cdot 2^n = 2^{n-1}$.
\qed

\textbf{Why this is weaker.}
The up-set $U$ is constructed from the internal structure of $\F$ and
need not be a star $[x,[n]]$.
To prove the full conjecture one would need $U = [x, [n]]$ for some specific
element $x$, which requires controlling not just the size of $U$ but its shape.
In particular, Remark 3.8 of \cite{nagel2023} notes that the $n$ star-upsets
form a vanishingly small fraction of all up-sets of size $\le 2^{n-1}$,
so the gap between the weakening and the conjecture is genuinely large.

\subsection{Averaging arguments and their limitations}

A natural first attempt is an \emph{averaging argument}: show that the average
element belongs to at least $\frac{1}{2}|\F|$ sets.
Formally, if we pick $x \in [n]$ uniformly at random, the expected count is
$\frac{1}{n}\sum_{x=1}^{n} |\{F \in \F : x \in F\}| = \frac{1}{n}\sum_{F \in \F}|F|$.
This is $\ge \frac{1}{2}|\F|$ iff the average set size is $\ge n/2$, which fails for
families with many small sets (e.g.\ $\F = \{\{1\}, \{1,2\}, \{1,3\}, \{1,2,3\}\}$ has
average set size $10/4 = 2.5 < 3/2 = n/2$ for $n = 3$... actually $n/2 = 1.5$,
so this doesn't give a counterexample; but $\F = \{\{1,2\},\{1,3\},\{2,3\},\{1,2,3\}\}$
has average size $10/4 = 2.5 < 2 = n/2$).
In any case, the conjecture is strictly weaker than the averaged version,
so the averaged version can fail even when the conjecture holds.

Gowers \cite{gowers2016} notes that a \emph{weighted} average might succeed:
does there exist a probability distribution $\pi$ on $[n]$ such that
$\mathbb{E}_{x \sim \pi}|\{F \in \F : x \in F\}| \ge \frac{1}{2}|\F|$?
This is trivially true if the conjecture holds (put all mass on the certifying element),
so the question must impose conditions on $\pi$ to avoid circularity.
One might ask for $\pi$ to be defined from $\F$ in a natural way, perhaps via
entropy maximisation.

Gowers also discusses a probabilistic construction of a potential counterexample:
choose each $k$-element subset of $[n]$ independently with probability $p$,
take $\F$ to be all unions of chosen sets.
If $k < n/2$ then typical unions have size $> n/2$ while the base sets
each have size $k$; one might hope the small sets are sparse enough
that every element is in few sets.
The construction fails: a back-of-envelope calculation shows no choice
of $k$ and $p$ makes the expected number of small sets dominate.
The conclusion is that \emph{random} union-closed families overwhelmingly
satisfy the conjecture, which suggests the hard cases are highly structured
rather than typical --- matching the tight-case structure we found computationally.

\subsection{Computational exploration}
\label{sec:comp-rl}

Beyond brute-force enumeration (Section~\ref{sec:enum}), we implemented
a proof-search framework based on the injection paradigm.

\textbf{Setup.}
For each sampled union-closed family we try to find a B-selection rule ---
a prescription for choosing $B(A) \in \F^+$ for each $A \in \F^-$ ---
that makes $\varphi(A) = A \cup B(A)$ injective.
We tested seven rules including two known strategies (K1, K2) and five novel ones:
{\small\texttt{union\_closest}} (pick $B \in \F^+$ minimising $|A \triangle B|$),
{\small\texttt{union\_superset\_pref}} (prefer $B \supseteq A$), and others.

\textbf{Results.}
Across 931 episodes with $n \in \{2, 3, 4, 5\}$, an LLM policy successfully
found injective maps in 84.8\% of episodes (excluding fallbacks).
Among the novel rules,
{\small\texttt{union\_closest}} achieved a 96.2\% success rate and
{\small\texttt{union\_superset\_pref}} achieved 91.6\%;
the structurally blunt {\small\texttt{union\_max}} succeeded only 40\% of the time.

\textbf{Significance.}
No single rule succeeded on every family.
In particular, the tight cases $\{\emptyset, A, [n]\setminus A, [n]\}$ require
careful element and rule selection, since both elements have equal frequency.
If a rule could be proved to yield an injective map for \emph{every}
nontrivial union-closed family, that would constitute a proof of the conjecture.
The computational evidence suggests {\small\texttt{union\_closest}} is a strong
candidate for further theoretical analysis.

% ---------------------------------------------------------------
\section{Progress and Obstacles}
% ---------------------------------------------------------------

\subsection{Computational verification for small \texorpdfstring{$n$}{n}}
\label{sec:enum}

We wrote a brute-force enumerator in Python that iterates over all
subsets of $\Pow(n)$, filters for nontrivial union-closed families,
and checks whether the conjecture holds for each one.
The results are summarised in Table~\ref{tab:enum}.

\begin{table}[h]
\centering
\begin{tabular}{rrrrc}
\hline
$n$ & Candidates & UC families & Tight cases & Conjecture \\
\hline
2 & 8 & 8 & 2 & \checkmark \\
3 & 128 & 90 & 5 & \checkmark \\
4 & 32{,}768 & 4{,}542 & 15 & \checkmark \\
\hline
\end{tabular}
\caption{Enumeration results. ``Tight'' means the best-element fraction
equals exactly $1/2$. Conjecture verified for all 4,640 families found.}
\label{tab:enum}
\end{table}

A \emph{tight} family is one where the best element sits at exactly
$\frac{1}{2}$; these are the hardest cases.
Inspecting the output, every tight family is either a complementary-pair
family $\{\emptyset, A, [n]\setminus A, [n]\}$ for some $A \subsetneq [n]$,
or the full power set $\Pow(n)$.
Both types are maximally symmetric --- no element is distinguishable
from any other by its frequency --- which matches the known extremal
construction from Nagel~\cite{nagel2023}.

\subsection{The scalability wall}

The search space for brute-force enumeration is $2^{|\Pow(n)|-1} = 2^{2^n - 1}$
candidate families, which is \emph{doubly exponential} in $n$.
Table~\ref{tab:timing} shows timed runs for $n \le 4$ and sampling-based
estimates for $n \ge 5$.

\begin{table}[h]
\centering
\begin{tabular}{rrrl}
\hline
$n$ & Candidates & Runtime & Feasible? \\
\hline
2 & $8$ & $< 1$ ms & Yes \\
3 & $128$ & $< 1$ ms & Yes \\
4 & $32{,}768$ & $72$ ms & Yes \\
5 & $\approx 2.1 \times 10^{9}$ & $\approx 6.2$ hrs & Borderline \\
6 & $\approx 9.2 \times 10^{18}$ & $\approx 5 \times 10^{6}$ yrs & No \\
7 & $\approx 1.7 \times 10^{38}$ & $\approx 10^{26}$ yrs & No \\
\hline
\end{tabular}
\caption{Candidate counts and wall-clock runtimes. Estimates for $n \ge 5$
obtained by sampling 2,000 random candidates and extrapolating.}
\label{tab:timing}
\end{table}

The key observation is that going from $n=5$ to $n=6$ does not merely
double the work --- it \emph{squares} it, since the exponent itself doubles.
Verification beyond $n=4$ would require fundamentally different approaches
such as incremental family construction, symmetry reduction under the
action of $S_n$, or SAT-based search.

% ---------------------------------------------------------------
\begin{thebibliography}{9}

\bibitem{bruhn2015}
H.~Bruhn and O.~Schaudt,
\textit{The Journey of the Union-Closed Sets Conjecture},
Graphs and Combinatorics \textbf{31}, pp.~2043--2074 (2015).

\bibitem{nagel2023}
N.~Nagel,
\textit{Notes on the Union Closed Sets Conjecture},
arXiv:2208.03803v2 (2023).

\bibitem{gowers2016}
W.T.~Gowers et al.,
\textit{Frankl's union-closed conjecture --- a possible Polymath project?},
\url{https://gowers.wordpress.com/2016/01/21/frankls-union-closed-conjecture-a-possible-polymath-project/} (2016).

\bibitem{balla2011}
I.~Balla,
\textit{Minimum density of union-closed families},
arXiv:1106.0369 (2011).

\bibitem{reimer2003}
D.~Reimer,
\textit{An Average Set Size Theorem},
Combinatorics, Probability and Computing \textbf{12}, pp.~89--93 (2003).

\bibitem{wojcik1999}
P.~W\'{o}jcik,
\textit{Union-closed families of sets},
Discrete Mathematics \textbf{199}, pp.~173--182 (1999).

\end{thebibliography}

\end{document}
